\chapter[Multiplicity Matrices]{Multiplicity Matrices}\label{ch:sec:multiplicity}
We now pass from an encoder to the integer data that will carry the representation.
\section{Definition and elementary identities}
\begin{definition}[Multiplicity matrix]\label{ch:def:multiplicity}
Let $E:\Mset\times\Kset\to\Cset$ be a deterministic Shannon encoder. Its \emph{multiplicity matrix} is the matrix $N=(n_{mc})_{m\in\Mset,c\in\Cset}$ defined by
$n_{mc}=\bigl|\{k\in\Kset:E(m,k)=c\}\bigr|.$
\end{definition}
\begin{lemma}\label{ch:lem:row-sum}
For every $m\in\Mset$,
\[
 \sum_{c\in\Cset}n_{mc}=|\Kset|,
 \qquad
 \Pr[C=c\mid M=m]=\frac{n_{mc}}{|\Kset|}.
\]
Moreover, $m\sim c$ if and only if $n_{mc}>0$.
\end{lemma}
\begin{proof}
For fixed $m$, every key has exactly one image under $k\mapsto E(m,k)$, so the sets counted by the entries in the $m$th row partition $\Kset$. Dividing by $|\Kset|$ gives the transition probability. The last assertion follows from Definition~\ref{ch:def:cloak} and the full support of $\pi$.
\end{proof}

\begin{lemma}[Posterior formula]\label{ch:lem:posterior-formula}
For every active ciphertext $c$ and $m\in\Mset$,
\[
 \Pr[M=m\mid C=c]
 =\frac{\pi(m)n_{mc}}{\sum_{m'\in\Mset}\pi(m')n_{m'c}}.
\]
\end{lemma}
\begin{proof}
Substitute the second identity of Lemma~\ref{ch:lem:row-sum} into Bayes' formula~\eqref{ch:eq:bayes}. The factor $|\Kset|^{-1}$ cancels from numerator and denominator.
\end{proof}
\section{The factorization criterion}
\begin{proposition}\label{ch:prop:factorization-criterion}
For an active ciphertext $c$, the posterior distribution on $\cloak(c)$ is uniform if and only if there exists a number $\beta(c)>0$ such that
\[
 \pi(m)n_{mc}=\beta(c)
 \qquad(m\in\cloak(c)).
\]
In that case,
\[
 \beta(c)=\frac{1}{|\cloak(c)|}\sum_{m\in\Mset}\pi(m)n_{mc}
 =\frac{|\Kset|\Pr[C=c]}{|\cloak(c)|}.
\]
\end{proposition}
\begin{proof}
By Lemma~\ref{ch:lem:posterior-formula}, the posterior is uniform on $\cloak(c)$ precisely when the numerator $\pi(m)n_{mc}$ is independent of $m$ on that set. Summing the resulting identity over the cloak gives the first formula for $\beta(c)$. The second follows from the law of total probability.
\end{proof}

\begin{corollary}\label{ch:cor:uo-factorization}
An encoder is universally optimal relative to $\pi$ if and only if there is a function $\beta$ on the active ciphertexts such that
\[
 \pi(m)n_{mc}=\beta(c)
 \qquad\text{whenever }n_{mc}>0.
\]
\end{corollary}
The corollary is the precise point at which a posterior condition becomes a combinatorial identity.
